%=============================================================================
% SECTION XVII.F: QUANTUM SUPERPOSITION AND THE ψ-FIELD
% Draft subsection for section_quantum.tex
% Closes the Page–Geilker gap; supersedes §XIV treatment in section_openproblems.tex
%=============================================================================

%-----------------------------------------------------------------------------
\subsection{Quantum Superposition and the $\psi$-Field: The Branch-by-Branch Picture}\label{sec:quantum-superposition}
%-----------------------------------------------------------------------------

\subsubsection{The Problem}

When matter is placed in spatial quantum superposition, what does $\psi$ do?  This question confronts any theory in which a classical field mediates gravitational effects.  In GR the question is acute: a mass in superposition at locations $A$ and $B$ appears to require spacetime that ``curves both ways,'' leading to the Penrose paradox and the broader problem of semiclassical gravity.

In DFD, there is one flat $\mathbb{R}^3$ with one scalar field $\psi$.  The field equation (static limit) is
\begin{equation}\label{eq:psi-source-static}
\nabla \cdot \bigl[\mu(|\nabla\psi|/a_\star)\,\nabla\psi\bigr] = -\frac{8\pi G}{c^2}\,\rho\,.
\end{equation}
The question reduces to: what is $\rho$ when the matter source is in a quantum superposition?

\subsubsection{Two Candidate Answers}

There are two logically distinct possibilities:

\paragraph{(i) Expectation-value sourcing.}
The field $\psi$ responds to the quantum expectation value of the mass density:
\begin{equation}\label{eq:expect-source}
\rho \;\to\; \langle\hat\rho\rangle = |c_A|^2\,\rho_A + |c_B|^2\,\rho_B\,.
\end{equation}
This is the analog of semiclassical gravity ($G_{\mu\nu} = 8\pi G\,\langle T_{\mu\nu}\rangle$).  It gives a single $\psi$ field for the system.  An earlier discussion (\S\ref{sec:open-penrose}) adopted this picture.

\paragraph{(ii) Branch-by-branch sourcing.}
Each branch of the matter superposition determines its own $\psi$ configuration, and the total quantum state is an entangled superposition of (matter, $\psi$) pairs:
\begin{equation}\label{eq:branch-state}
|\Psi_{\rm total}\rangle = c_A\,|A,\,\psi_A\rangle + c_B\,|B,\,\psi_B\rangle\,,
\end{equation}
where $\psi_A$ is the unique solution of Eq.~\eqref{eq:psi-source-static} with source $\rho_A$, and $\psi_B$ the unique solution with source $\rho_B$.

These two pictures are \emph{not} equivalent.  Expectation-value sourcing is inconsistent with standard quantum mechanics in general (the Page--Geilker argument~\cite{PageGeilker1981}; see also~\cite{EppleyHannah1977,Carlip2008}).  The branch-by-branch picture preserves full quantum linearity.  We now argue that DFD, in the quasi-static limit relevant to all laboratory experiments, naturally yields the branch-by-branch picture.

\subsubsection{The Quasi-Static Argument}

The core of the argument rests on three facts:

\begin{enumerate}
\item \textbf{Ellipticity in the static limit.}  In the time-independent (or quasi-static) limit, the DFD field equation~\eqref{eq:psi-source-static} is a quasilinear \emph{elliptic} PDE.  No propagating degrees of freedom exist: $\psi$ is entirely determined by the instantaneous source distribution $\rho(\mathbf{x})$ and boundary conditions at infinity (Sec.~\ref{sec:elliptic}).

\item \textbf{Uniqueness of elliptic solutions.}  For a given source $\rho$, the monotonicity condition (A4) on $\mu$ (Sec.~\ref{sec:elliptic}) ensures that the boundary value problem has a \emph{unique} solution $\psi[\rho]$.  The solution is a deterministic functional of the source.

\item \textbf{Superposition of quantum branches.}  Standard quantum mechanics requires that if the matter Hilbert space admits a superposition $c_A|A\rangle + c_B|B\rangle$, then the total state of (matter $+$ everything that couples to matter) is a corresponding superposition.  Since $\psi[\rho_A]$ and $\psi[\rho_B]$ are each uniquely determined, the combined state is Eq.~\eqref{eq:branch-state}.
\end{enumerate}

The logic is:
\begin{quote}
\emph{Elliptic equation $+$ unique solution for each source $\Rightarrow$ each matter branch gets its own $\psi$ $\Rightarrow$ the quantum state is a superposition of (matter, $\psi$) pairs.}
\end{quote}

This is precisely analogous to the Coulomb field in nonrelativistic QED: the instantaneous $1/r$ potential is not an independent degree of freedom but a constraint determined by the charge distribution.  When a charge is in superposition, the Coulomb field goes branch-by-branch.  No one worries about this because the Coulomb potential satisfies Poisson's equation---an elliptic PDE with unique solutions.  The DFD $\psi$ field inherits this same structure in the quasi-static limit.

\subsubsection{Validity of the Quasi-Static Limit}

The quasi-static approximation holds when the $\psi$-field equilibration time $\tau_\psi$ is much shorter than the quantum coherence time $\tau_{\rm coh}$ of the matter system.  For perturbations about a smooth background, the DFD field equation has a characteristic propagation speed $c_\psi \leq c$ (Sec.~\ref{sec:open-hyperbolicity}).  Thus:
\begin{equation}\label{eq:quasi-static-criterion}
\tau_\psi \sim \frac{L}{c_\psi}\,, \qquad \text{quasi-static limit valid when}\quad \frac{L}{c_\psi} \ll \tau_{\rm coh}\,,
\end{equation}
where $L$ is the spatial extent of the matter superposition.

For all currently conceivable laboratory experiments:
\begin{itemize}
\item \textbf{Atom interferometers:} $L \sim 1\,\text{m}$, so $\tau_\psi \lesssim 3\,\text{ns}$.  Coherence times are $\tau_{\rm coh} \sim 1\,\text{s}$.  Ratio: $\tau_\psi/\tau_{\rm coh} \lesssim 10^{-9}$.
\item \textbf{Macroscopic superposition proposals} (e.g., $10^{-14}\,\text{kg}$ oscillators): $L \sim 10^{-6}\,\text{m}$, $\tau_\psi \lesssim 10^{-14}\,\text{s}$, $\tau_{\rm coh} \sim 1\,\text{s}$.  Ratio: $\lesssim 10^{-14}$.
\item \textbf{MAQRO-class experiments:} $L \sim 10^{-7}\,\text{m}$, $\tau_\psi \lesssim 10^{-15}\,\text{s}$, $\tau_{\rm coh} \sim 100\,\text{s}$.  Ratio: $\lesssim 10^{-17}$.
\end{itemize}

In every case, $\psi$ equilibrates to its constraint-surface value essentially instantaneously on the timescale of quantum evolution.  The quasi-static limit is valid by many orders of magnitude.

\subsubsection{The Temporal Sector: An Honest Assessment}\label{sec:temporal-honest}

The full DFD field equation is not purely elliptic.  Time-dependent perturbations about smooth backgrounds obey a \emph{hyperbolic} wave equation (Sec.~\ref{sec:open-hyperbolicity}), with propagation speed determined by the perturbation metric $\mathcal{G}^{\mu\nu}$.  In the hyperbolic sector, $\psi$ carries genuine propagating degrees of freedom---it has its own dynamics beyond the constraint structure.

This means:

\begin{enumerate}
\item The branch-by-branch argument above is \emph{rigorous} in the quasi-static limit (where the hyperbolic sector is frozen out).

\item For processes where $\psi$-propagation is dynamically relevant ($\tau_\psi \gtrsim \tau_{\rm coh}$), the correct quantum treatment of $\psi$ becomes a nontrivial open question.  In this regime, one must either:
\begin{itemize}
\item[(a)] Quantize the $\psi$-field perturbations, yielding a quantum field theory of $\psi$ coupled to matter (the full ``quantum gravity'' problem within DFD), or
\item[(b)] Accept that the quasi-static branch-by-branch picture receives small corrections of order $\tau_\psi/\tau_{\rm coh}$.
\end{itemize}

\item For all laboratory quantum experiments, $\tau_\psi/\tau_{\rm coh} < 10^{-9}$, and option (b) gives corrections that are experimentally undetectable.  The full quantum field theory of $\psi$ is therefore an open theoretical question but not an experimental concern in any foreseeable experiment.
\end{enumerate}

\paragraph{Action scaling.}
An additional reason why the hyperbolic sector is irrelevant for laboratory quantum coherence: the action of the $\psi$ field scales as $S_\psi \sim (M_P/a_\star)^2 \gg \hbar$ (Sec.~\ref{sec:open-uv}).  Quantum fluctuations of $\psi$ are suppressed by this enormous action ratio.  The field is effectively classical; only its \emph{sourcing} enters quantum mechanics, and the branch-by-branch picture handles this correctly.

\subsubsection{Zero Gravitational Decoherence: A Sharp Prediction}\label{sec:zero-decoherence}

The branch-by-branch picture yields a concrete, falsifiable prediction:

\begin{tcolorbox}[colback=green!5!white, colframe=green!60!black, title=DFD Prediction: Zero Gravitational Decoherence]
In the quasi-static limit, the $\psi$-field introduces \textbf{no intrinsic decoherence} to quantum superpositions.  Massive particles in spatial superposition evolve according to standard unitary quantum mechanics.  Any observed decoherence must be attributed to environmental (nongravitational) sources.
\end{tcolorbox}

This prediction sharply contrasts with two prominent alternatives:

\paragraph{Di\'osi--Penrose (DP) mechanism.}
The DP proposal~\cite{Diosi1989,Penrose1996} predicts that spatial superpositions of masses decohere at a rate
\begin{equation}\label{eq:dp-rate}
\Gamma_{\rm DP} = \frac{1}{\hbar}\int d^3x\,d^3x'\;\frac{G\,\Delta\rho(\mathbf{x})\,\Delta\rho(\mathbf{x}')}{|\mathbf{x} - \mathbf{x}'|}\,,
\end{equation}
where $\Delta\rho = \rho_A - \rho_B$ is the mass density difference between branches.  The physical basis is that ``spacetime cannot support two incompatible geometries,'' so gravity forces collapse.

In DFD, there is no curved spacetime---only flat $\mathbb{R}^3$ with a scalar field.  The scalar field $\psi$ supports branch-by-branch solutions without any inconsistency.  There is no mechanism to generate $\Gamma_{\rm DP}$.

\paragraph{Continuous Spontaneous Localization (CSL).}
CSL models~\cite{Bassi2013} postulate a universal noise field that collapses superpositions, with the rate scaling with mass.  DFD has no such noise field.  The $\psi$-medium is a deterministic classical background (in the quasi-static limit); it does not inject stochasticity into quantum evolution.

\paragraph{Experimental discrimination.}
Current and planned experiments can distinguish these predictions:
\begin{itemize}
\item \textbf{MAQRO}~\cite{MAQRO2022}: Space-based matter-wave interferometry with $10^{9}$--$10^{10}$\,amu particles.  DP predicts anomalous decoherence scaling as $M^2$; DFD predicts standard quantum behavior.
\item \textbf{Macroscopic oscillator experiments}: Proposals to place $\sim$$10^{-14}\,\text{kg}$ objects in superposition~\cite{Bouwmeester2003}.  DP predicts decoherence times $\sim 10^{-7}\,\text{s}$; DFD predicts no gravitational decoherence.
\item \textbf{Atom interferometers in drop towers}: Extended free-fall times increase $\tau_{\rm coh}$ without environmental decoherence, testing whether anomalous decoherence appears.
\end{itemize}

\subsubsection{Relation to the Expectation-Value Treatment}\label{sec:expect-value-limit}

The expectation-value sourcing described in Sec.~\ref{sec:open-penrose} is not wrong---it is the \emph{decoherent limit} of the branch-by-branch picture.  The connection is straightforward:

When environmental decoherence has fully decohered the matter superposition (so that the density matrix is diagonal in the position basis), the reduced density matrix for $\psi$ is
\begin{equation}
\hat\rho_\psi = |c_A|^2\,|\psi_A\rangle\langle\psi_A| + |c_B|^2\,|\psi_B\rangle\langle\psi_B|\,.
\end{equation}
In this classical mixture, the expectation value of any $\psi$-dependent observable $\hat{\mathcal{O}}$ is
\begin{equation}
\langle\hat{\mathcal{O}}\rangle = |c_A|^2\,\mathcal{O}[\psi_A] + |c_B|^2\,\mathcal{O}[\psi_B]\,.
\end{equation}
For \emph{linear} functionals of $\psi$ (or when $\psi_A \approx \psi_B$), this reduces to $\mathcal{O}[\langle\psi\rangle]$---the expectation-value result.

The expectation-value treatment is therefore valid in two limits:
\begin{enumerate}
\item When the matter has fully decohered (classical limit);
\item When the $\psi$-field difference between branches is small compared to the background (linearized limit).
\end{enumerate}
Both limits apply to all astrophysical and cosmological contexts.  The branch-by-branch picture is necessary only for laboratory quantum superposition experiments.

The treatment in Sec.~\ref{sec:open-penrose} is accordingly superseded by the present analysis for quantum-coherent matter.  The earlier expectation-value statements remain correct for decohered (classical) matter.

\subsubsection{Page--Geilker Consistency}\label{sec:page-geilker}

Page and Geilker~\cite{PageGeilker1981} argued that semiclassical gravity ($G_{\mu\nu} = 8\pi G\,\langle T_{\mu\nu}\rangle$) makes predictions inconsistent with observation: a Cavendish-balance experiment driven by a quantum random number generator shows that gravity tracks the \emph{actual} mass configuration, not the expectation value.

DFD is fully consistent with Page--Geilker because:
\begin{enumerate}
\item In the branch-by-branch picture, each measurement outcome selects a branch $|A,\psi_A\rangle$ or $|B,\psi_B\rangle$.
\item After measurement, the observer is in a definite branch with a definite $\psi$.
\item No observer ever sees $\psi$ responding to the expectation value of a superposition that has been collapsed (or branched) by measurement.
\end{enumerate}

This is the same resolution that applies to the electromagnetic field: after measuring which slit a photon passed through, the observer sees the Coulomb field of the photon at its actual position, not the expectation value over both slits.

\subsubsection{The BMV Experiment: A Testable Prediction}\label{sec:bmv-prediction}

The branch-by-branch picture makes a direct prediction for the Bose--Marletto--Vedral (BMV) experiment~\cite{Bose2017,MarlettoVedral2017}.  In the BMV protocol, two masses $m_1$ and $m_2$ are each placed in spatial superposition, and gravitational interaction between them is tested for the ability to generate entanglement.

\paragraph{Setup.}
Each mass is in a superposition of two locations separated by $\Delta x$.  In the branch-by-branch picture, the joint state evolves as:
\begin{equation}
|L_1, L_2, \psi_{LL}\rangle + |L_1, R_2, \psi_{LR}\rangle + |R_1, L_2, \psi_{RL}\rangle + |R_1, R_2, \psi_{RR}\rangle\,,
\end{equation}
where $\psi_{ij}$ is the unique solution of Eq.~\eqref{eq:psi-source-static} for the source configuration with mass~1 at position $i$ and mass~2 at position $j$.

\paragraph{Entanglement generation.}
Because $\psi_{LL} + \psi_{RR} \neq \psi_{LR} + \psi_{RL}$ (the $\psi$-field equation is nonlinear in general, and even in the linear regime the gravitational phase accumulated differs between branches), the four branches acquire different gravitational phases.  The resulting state is generically \emph{entangled}: it cannot be written as a product state of mass~1 and mass~2.

\paragraph{DFD prediction.}
DFD predicts that the BMV experiment \emph{will} observe gravitationally induced entanglement at the rate predicted by Newtonian gravity (since the quasi-static limit applies with enormous margin for laboratory-scale masses and separations).  The entanglement phase is:
\begin{equation}\label{eq:bmv-phase}
\Delta\phi_{\rm grav} \approx \frac{Gm_1 m_2}{\hbar}\left(\frac{1}{d - \Delta x} - \frac{1}{d}\right)\,T\,,
\end{equation}
where $d$ is the mean separation and $T$ is the interaction time.  This is the same prediction as quantized linearized gravity, because in the quasi-static regime the $\psi$-field is a constraint, not an independent quantum degree of freedom.

\paragraph{Contrast with semiclassical gravity.}
Semiclassical gravity (expectation-value sourcing) predicts \emph{no} gravitationally induced entanglement, because the gravitational field responds only to the mean, washing out the which-path information needed for entanglement.  A positive BMV result would therefore:
\begin{itemize}
\item Rule out expectation-value sourcing (semiclassical gravity);
\item Be \emph{consistent} with DFD's branch-by-branch picture;
\item Not, by itself, prove that gravity is ``quantized'' in the particle-physics sense---only that the gravitational field tracks branches rather than expectation values.
\end{itemize}

\subsubsection{Summary of Claim Status}

\begin{table}[h]
\centering
\caption{Quantum superposition of $\psi$: claim status.}\label{tab:quantum-psi-status}
\begin{ruledtabular}
\begin{tabular}{lll}
\textbf{Claim} & \textbf{Status} & \textbf{Basis} \\
\hline
Branch-by-branch in quasi-static limit & \textbf{Theorem} & Elliptic uniqueness \\
Quasi-static valid for all lab experiments & \textbf{Established} & $\tau_\psi/\tau_{\rm coh} < 10^{-9}$ \\
Zero gravitational decoherence & \textbf{Prediction} & Falsifiable by MAQRO \\
BMV entanglement positive & \textbf{Prediction} & Quasi-static + branch picture \\
Full QFT of $\psi$ in hyperbolic sector & \textbf{Open} & Needed if $\tau_\psi \sim \tau_{\rm coh}$ \\
$\langle\rho\rangle$ sourcing as decoherent limit & \textbf{Derived} & Reduced density matrix \\
Page--Geilker consistency & \textbf{Established} & Branch selection at measurement \\
\end{tabular}
\end{ruledtabular}
\end{table}

\begin{tcolorbox}[colback=blue!5!white, colframe=blue!60!black, title=Branch-by-Branch Picture: Summary]
\small
In the quasi-static limit (valid for all laboratory experiments by $\geq 9$ orders of magnitude), the DFD field $\psi$ goes branch-by-branch with quantum superpositions of matter.  This follows from elliptic uniqueness: each matter configuration determines a unique $\psi$, so superposed matter yields superposed (matter, $\psi$) pairs.

\textbf{Consequences:}
\begin{enumerate}[nosep, leftmargin=*]
\item No gravitational decoherence (sharp contrast with Di\'osi--Penrose)
\item Page--Geilker consistency (no conflict with Cavendish experiments)
\item BMV entanglement predicted (contrast with semiclassical gravity)
\item Expectation-value sourcing recovered as the decoherent limit
\end{enumerate}

\textbf{Open:} The full quantum field theory of $\psi$ in the hyperbolic (dynamical) sector.  This is a theoretical question, not an experimental concern at accessible energy scales.
\end{tcolorbox}
